\documentclass{exam}
\usepackage{graphicx} % Required for inserting images
\usepackage{fancyhdr}
\usepackage{amssymb}
\usepackage{amsmath}
\usepackage{mathtools}
\usepackage[most]{tcolorbox}
\usepackage[dvipsnames]{xcolor}
\usepackage{blindtext}
\setlength{\parindent}{0pt}
% if you ever want to enter images use the following code%
%\begin{center}%
    %\includegraphics[width=0.8\textwidth]{Screenshot (22).png}%
%\end{center}%

\newtheorem{thm}{Theorem}[section]
\newtheorem{lem}{Lemma}[section]
\newtheorem{defn}{Definition}[section]
\newtheorem{cor}{Corollary}[section]
\newtheorem{axm}{Axiom}[section]

\tcolorboxenvironment{thm}{
  colback=lime,
  colframe=black,
  breakable
}

\tcolorboxenvironment{lem}{
  colback=yellow,
  colframe=black,
  breakable
}

\tcolorboxenvironment{defn}{
  colback=pink,
  colframe=black,
  breakable
}

\tcolorboxenvironment{cor}{
  colback=SpringGreen,
  colframe=black,
  breakable
}

\tcolorboxenvironment{axm}{
  colback=white,
  colframe=black,
  breakable
}


\fancyhf{}
\fancyhead[R]{   Omkar Tasgaonkar\\UID: 206-624-454}
\renewcommand\headrulewidth{0.5pt}
\pagestyle{fancy}
\author{Omkar Tasgaonkar}
\setlength{\headsep}{0.5in}
\fancyfoot[C]{(\thepage)}
\renewcommand\footrulewidth{0.5pt}

\fancypagestyle{firstpage}{
    \fancyhf{}
    \fancyhead[C]{\textbf{Notes on ROS}}
    \fancyhead[R]{Omkar Tasgaonkar\\UID: 206-624-454}
    \renewcommand\headrulewidth{0.5pt}
    \fancyfoot[C]{(\thepage)}
    \renewcommand\footrulewidth{0.5pt}
}



\begin{document}
\thispagestyle{firstpage}

\section{Nodes}
The idea of nodes is that each system within a robot is a node in and of itself; with nodes being able to send and receive messages from other nodes. To find a list of active nodes, we use the command:
\begin{verbatim}
$ros2 node list
\end{verbatim}
So if we run the following command to open turtlesim, we get the following output:
\begin{verbatim}
Input: 
$ros2 run turtlesim turtlesim_node
$ros2 node list

Output:
/turtlesim
\end{verbatim}

\subsection{Remapping}
We can change values that each node possesses via remapping commands. For example, the following command changes our turtlesim node's name:
\begin{verbatim}
$ros2 run turtlesim turtlesim_node --ros-args --remap__node:=my_turtle
\end{verbatim}
In this case, we run the turtlesim once more, creating a new sim instance; however, the node is now named my\_turtle.

\section{Topics}
The fundamental idea behind topics is the publisher-subscriber model. The publishers are the nodes who send the message and the subscribers receive it. Multiple publishers and subscribers can connect to one topic at the same time.\\

A message sent from one publisher is transmitted to all subscribers. 
\end{document}
